\subsection{Finite puncture sets and their class-two quotients}

\begin{result}{E.40}{Proved}{The simultaneous radical lemma used for finite puncture sets}
Let $4\mid m$, and let $d_1,\ldots,d_s$ be positive rational numbers.
Suppose each $d_i$ has a private rational prime $P_i$ with
$v_{P_i}(d_i)=\pm1$ and $v_{P_i}(d_j)=0$ for $j\ne i$.
Assume that $F=\mathbb Q(\mu_{2m})$ contains all $\sqrt{d_i}$.
Then, for the Galois extension
\[
 L=F(d_1^{1/m},\ldots,d_s^{1/m}),
 \qquad \operatorname{Gal}(L/F)\simeq(\mathbb Z/(m/2))^s,
\]
the maximal abelian subfield of $L/\mathbb Q$ is exactly $F$.
For every multiple $M$ of $2m$, every unit residue class modulo $M$
which is one modulo $2m$ contains arbitrarily large primes splitting
completely in $L$. This confirms SECTIONS\_OVER\_Q A.47, the
arithmetic input to its A.48.
\end{result}
\begin{proof}
Put $h=m/2$ and choose the positive real roots $\alpha_i=d_i^{1/m}$.
If a positive rational number $d$ has an $m$th root in $F$, its
positive root $\alpha$ belongs to $F$, since $F$ contains $\mu_m$.
The real field $\mathbb Q(\alpha)$ is Galois, because it is a
subfield of an abelian Galois extension. Its conjugates can only be
the real roots $\alpha,-\alpha$ of $X^m-d$. Hence
$\alpha^2\in\mathbb Q$ and $d$ is an $h$th power in $\mathbb Q$.
Apply this to $d=\prod_i d_i^{e_i}$. The private valuations force
$h\mid e_i$ whenever this product is an $m$th power in $F$.
Conversely $d_i^h=(\sqrt{d_i})^m$. Thus these classes generate
$(\mathbb Z/h)^s$ in $F^\times/(F^\times)^m$.
The perfect Kummer pairing identifies its dual with
$N=\operatorname{Gal}(L/F)$; its precise source is
\citet[Theorem 5.30 and Remark 5.32, pp.~75--76]{RootMilneFT}.
Use coordinates in which a basis element multiplies one $\alpha_i$
by $\zeta_m^2$ and fixes the others.

Let $G=\operatorname{Gal}(L/\mathbb Q)$ and let $c$ be complex
conjugation. It fixes every $\alpha_i$ and acts by $-1$ on $N$.
Consequently $2N\subseteq[G,G]\subseteq N$. If
$g(\alpha_i)=\zeta_m^{k_i}\alpha_i$, direct composition gives
\[
 [c,g](\alpha_i)=\zeta_m^{-2k_i}\alpha_i.
\]
The parity vector $(k_i\bmod2)$ is the action on the independent
quadratic roots $\sqrt{d_i}$. Independence again follows from the
private valuations. Restriction from $G$ to their multiquadratic
field realizes every parity vector. The displayed commutators
therefore surject onto $N/2N$. Hence $[G,G]=N$ and $L^{\rm ab}=F$.

For $2m\mid M$, we have $L\cap\mathbb Q(\mu_M)=F$.
The identity automorphism of $L$ and the specified cyclotomic
residue automorphism agree on that intersection, so define an
actual element of the Galois group of the compositum. Its conjugates
retain both restrictions. Chebotarev, in the exact form of
\citet[Theorem 8.31, p.~147]{CohMilneANT}, supplies the primes.
\end{proof}

\begin{result}{E.41}{Proved}{SECTIONS\_OVER\_Q A.48 confirmed for every fixed finite puncture set}
Let $S\subset\mathbb Q$ be finite with $k=|S|\ge2$. There are
completed classes $U_c\in\widehat{\mathbb Q^\times}$ and actual
points $x_p\in\mathbb Q_p\setminus S$ at every finite prime such
that $(U_c)_p=\kappa_p(x_p-c)$ for every $c\in S$. Each $U_c$ has
infinite support, with infinitely many private support primes absent
from every other coordinate. There is also a real point with the
specified completed sign classes. In particular this local tuple is
not one rational point. The set $S$ is fixed before the construction;
no compatibility as $S$ varies is asserted.
\end{result}
\begin{proof}
We first treat integer punctures $c_1,\ldots,c_k$. Choose distinct
odd primes $P_i$ larger than every nonzero $|c_a-c_b|$, and by CRT
choose $R>\max_i c_i$ with $R\equiv c_i+P_i\pmod{P_i^2}$.
Set
\[
 a_{i,0}=R-c_i>0,\quad
 D_0=\prod_i a_{i,0}\prod_{i<t}|c_i-c_t|,\quad
 N_0=8\operatorname{rad}(D_0),\quad n_j=N_0(j+1)!.
\]
The $P_i$ are private primes of valuation one for these initial
coefficients and divide no puncture difference. Every quadratic
root $\sqrt{|(c_i-c_t)/a_{t,0}|}$ lies in
$\mathbb Q(\mu_{N_0})$. Indeed that cyclotomic field contains $i$,
$\sqrt2=\zeta_8+\zeta_8^{-1}$ and $\sqrt P$ for each odd
$P\mid D_0$, by the quadratic Gauss sum of square $(-1/P)P$.
Its norm-square computation is the one in E.102. Factoring the
rational numbers proves the assertion.

We give the point-count estimate used uniformly in the induction.
For $p\nmid m$, distinct residues $c_i$ and units $a_i$, the curve
\[
 T-c_i=a_iX_i^m\ (1\le i\le k),\qquad \prod_iX_i\ne0
\]
is a smooth geometrically connected cover of the punctured line of
degree $m^k$. Independence of the Kummer functions is witnessed by
their valuations at the distinct finite punctures; E.40's Kummer
pairing applies over the algebraically closed function field as
well. The ramification index is $m$ at each of the $k$ punctures
and at infinity. At infinity all functions have a simple pole, and
their ratios have roots in the strict local unit ring, so a single
$m$th root of its uniformizer accounts for all inertia. Thus
\[
 2g-2=(k-1)m^k-(k+1)m^{k-1},\qquad b=(k+1)m^{k-1}
\]
are the genus formula and the number of removed geometric points.
The Weil bound gives at least $p+1-2g\sqrt p-b$ interior points
\cite[(1.5.1), Theorem (1.6), pp.~275--276]{RootDeligne1974}.
Set $B_k(m)=((k+1)m^k+3)^2$. At
$t_0=(k+1)m^k+3$ we have
$t_0-2g=2m^k+(k+1)m^{k-1}+1>0$ and
$t_0(t_0-2g)+1-b>0$. The same quadratic is increasing thereafter.
Hence $p>B_k(m)$ guarantees an interior point, which lifts by the
unit derivatives of the equations in $X_i$.

Inductively we keep positive integer coefficients $a_{i,j}$, each
initial value times eighth powers of fresh primes. The private
valuations stay fixed. Put $V_j(p)=\max_i v_p(a_{i,j})$. Maintain
\[
 T_p-c_i=a_{i,j}X_{i,p}^{n_j},\qquad
 X_{i,p}\in\mathbb Z_p^\times \quad\hbox{for every finite }p,
\]
and, for $p\le B_k(n_j)$, maintain
\[
 v_p((a_{i,j}+c_i)-(a_{1,j}+c_1))\ge
 H_j(p):=V_j(p)+2v_p(2n_{j+1})+j+3.
\]
At stage zero $T_p=R$ and all $X_{i,p}=1$ work.
Write $n=n_j,m=n_{j+1},a_i=a_{i,j}$. Let $\mathcal T$ contain
all primes at most $B_k(m)$, all old coefficient primes and all
puncture-difference primes. For $p\in\mathcal T$ set
\[
 L_p=V_j(p)+2v_p(2n_{j+2})+j+4,
 \qquad M=\prod_{p\in\mathcal T}p^{L_p}.
\]
This contains every divisor of $2m$ and $2m\mid M$.
At a prime $p\mid2m$, choose the old point with
$T_p=a_1+c_1$ and $X_{1,p}=1$. For $i\ne1$ the quotient
$(T_p-c_i)/a_i$ differs from one with valuation at least
$2v_p(2m)+j+3$. Newton's estimate for the $n$th root, as proved
in E.102, gives a root with
$X_{i,p}\equiv1\pmod{p^{v_p(2m)}}$.
The needed old precision applies: such a prime either divides
$N_0$ or is at most $j+2$, so is below $B_k(n)$.
For the other primes of $\mathcal T$, choose any maintained point.

Choose new primes $q_1,\ldots,q_k$ successively. At the $i$th
choice put $b_t=a_tq_t^n$ for $t<i$ and $b_t=a_t$ for $t\ge i$.
For $t\ne i$ define $d_{it}=|(c_i-c_t)/b_t|$ and use the field
\[
 F=\mathbb Q(\mu_{2m}),\qquad L_i=F(d_{it}^{1/m}:t\ne i).
\]
Its radicands have private valuations $-1$ at the primes $P_t$;
all their squareclasses are fixed initial squareclasses. E.40
therefore proves $L_i^{\rm ab}=F$. Since an $m$th root of $-1$
lies in $F$, complete splitting also supplies roots of the signed
radicands.

Impose $q_i\equiv X_{i,p}\pmod{p^{L_p}}$ at all
$p\in\mathcal T$. At each previously chosen $q_h$, $h<i$, impose
\[
 q_i\equiv\rho_{hi}\pmod{q_h},\qquad
 \rho_{hi}^{n}\equiv(c_h-c_i)/a_i\pmod{q_h}.
\]
The residue $\rho_{hi}$ exists because $q_h$ was chosen to split
in the radical field containing an $m$th root of this number, and
$n\mid m$. The combined CRT class is a unit class equal to one
modulo $2m$. E.40 supplies an arbitrarily large prime $q_i$ in
that class splitting completely in $L_i$. We require it larger
than $B_k(m)$, every old or newly chosen prime and $n_{j+2}$;
thus all primes remain distinct and all required derivatives at
new primes are units. Define $a_{i,j+1}=a_iq_i^n$.

We check all new local conditions. At $p\in\mathcal T$,
$a_iq_i^n+c_i\equiv T_p\pmod{p^{L_p}}$ for every $i$.
Their differences satisfy the next displayed precision, since
$V_{j+1}(p)=V_j(p)$. Set $T'_p=a_{1,j+1}+c_1$.
Then $(T'_p-c_i)/a_{i,j+1}$ differs from one with valuation at
least $L_p-V_j(p)>2v_p(m)$, so has a unit $m$th root.
At a new prime $q_h$, set $T=c_h+a_hq_h^n$ and $X_h=1$.
For $i<h$, the other quotient reduces to $(c_h-c_i)/b_i$,
an $m$th power by the splitting imposed at $q_h$.
For $i>h$, it reduces to one by the subsequent CRT condition.
All are units, since $q_h$ divides no puncture difference, and
simple-root Hensel lifting applies. At every remaining prime the
coefficients and puncture differences are units, $m$ is invertible,
and $p>B_k(m)$; the proved point-count bound applies. This completes
the induction, including all dyadic and repeated-coefficient primes.

The completed products $U_i=a_{i,0}\prod_jq_{i,j}^{n_j}$ converge
because $n_j$ is eventually divisible by every integer. At any fixed
$p$ their ordinary valuations eventually stabilize. Their unit parts
converge in $\mathbb Q_p^\times$: each finite unit quotient is
annihilated by every sufficiently late $n_j$. The limits $y_{i,p}$
are nonzero, have the asserted completed Kummer classes, and the
precision $H_j(p)\to\infty$ gives $y_{i,p}+c_i=x_p$ independent
of $i$. Every $U_i$ has infinitely many private positive valuations,
so the tuple cannot be rational. All its real signs are positive;
choose $x_\infty>\max_i c_i$.

Finally choose an integer $D>0$ with $Dc_i\in\mathbb Z$ for
rational punctures. The $\mathbb Q$-isomorphism $t\mapsto Dt$
sends $t-c_i$ to $(T-Dc_i)/D$ in the indicated direction.
Apply the integer construction and divide the local points by $D$
and each completed class by $\kappa(D)$. This continuous translation
commutes with every localization, changes only finitely many valuation
coordinates, and preserves the required real signs. It proves the
statement over the original rational punctures.
\end{proof}

\begin{result}{E.42}{Proved}{SECTIONS\_OVER\_Q A.50 confirmed by pairwise forgetful maps}
For $V_S=\mathbb A^1_\mathbb Q\setminus\{c_1,\ldots,c_k\}$,
every completed tuple locally represented by $(x_p-c_i)_i$ at every
finite prime has a continuous section of its actual integral
class-two arithmetic quotient with that abelian tuple. The examples
of E.41 give neither rational-point nor cuspidal section classes.
The local central coordinates of these global sections are not
prescribed by this assertion.
\end{result}
\begin{proof}
Choose a rational base point $b\in V_S(\mathbb Q)$. The meridians
at the finite punctures identify the geometric abelian group with
$T=\widehat{\mathbb Z}(1)^k$ and the next central layer with
$M=\bigwedge^2T=\prod_{i<j}\widehat{\mathbb Z}(2)e_{ij}$.
These are continuous Galois-equivariant isomorphisms: the Kummer
degree of $t-c_i$ on the $j$th meridian is $\delta_{ij}$, and
commutators have the exterior-square action. The integral class-two
coordinate law is
\[
 (a,h)(a',h')=(a+a',h+h'+(a_i a'_j)_{i<j}).
\]
Word collection proves its universal property for every finite
class-two quotient, hence for the profinite group. It also gives
the continuous set section $a\mapsto(a,0)$.
Relative to the point-$b$ splitting the given abelian cocycle has
coordinates $U_i-\kappa(b-c_i)$. Pull back the next central
extension. Its actual continuous defect is a class in $H^2(G_\mathbb Q,M)$.

For $i<j$ the morphism
\[
 V_S\longrightarrow U,\qquad
 z=(t-c_i)/(c_j-c_i),\qquad 1-z=-(t-c_j)/(c_j-c_i)
\]
forgets all other punctures. It induces a continuous homomorphism
of arithmetic extensions, based at $b$ and its image. On $M$ its
map is the projection $e_{ij}\mapsto1$ and $e_{ab}\mapsto0$ for
other pairs. This follows on the surviving meridians; any path
change is inner and does not alter this central map.
The image abelian pair is
\[
 (U_i/\kappa(c_j-c_i),\ \kappa(-1)U_j/\kappa(c_j-c_i)),
\]
which is ordinary at every finite prime. A.25 proves its actual
continuous class-two obstruction is zero. Naturality of the defect
therefore makes every coordinate of the original obstruction zero.
Continuous cochains commute with this finite product of modules,
so the whole class is zero; a continuous primitive gives the section.

For a rational point each $t-c_i$ has finite support. At the cusp
$c_i$, every coordinate $t-c_j$, $j\ne i$, specializes to the
nonzero rational $c_i-c_j$ and again has finite support. At infinity,
with parameter $z=1/t$, all $t-c_i=z^{-1}(1-c_i z)$ have the same
completed Kummer class; their principal factors have compatible
roots with value one. The private support sets in E.41 exclude
all these possibilities already in the abelian quotient.
\end{proof}

\begin{result}{E.43}{Proved}{The exact ordinary-point scope of SECTIONS\_OVER\_Q A.196}
Let $k=\mathbb Q(\alpha)$ with $\alpha^3=3$, and put
$A=649^8\alpha$, $U=\alpha^2-2$, $B=U^3$.
At the unique place above three, the abelian pair modulo ninth powers
is neither an ordinary-point pair nor a cuspidal pair. Its reduction
modulo cubes is an ordinary-point pair. Thus the first failure along
these two nontrivial three-power Kummer truncations occurs modulo
nine. This confirms GT\_TOWER\_CHECK C.2.3 and the local scope
qualification of SECTIONS\_OVER\_Q A.196. It does not certify the
latter statement's separate global continuous nonvanishing assertion.
\end{result}
\begin{proof}
In $k_3=\mathbb Q_3(\alpha)$ use $v(\alpha)=1$, so $v(3)=3$.
The polynomial is Eisenstein, giving the unique place and residue
field $\mathbb F_3$. We have the exact identities
\[
 U-1=\alpha^2(1-\alpha),\qquad B-1=-6\alpha^2U,
 \qquad v(A)=1,\quad v(U-1)=2,\quad v(B-1)=5.
\]
For $Y\equiv1\pmod\alpha$, expansion of $(1+z)^3$ shows
$v(Y^3-1)\ge3$; another cubing gives $v(Y^9-1)\ge6$.
An ordinary pair with these ninth-power classes would give
\[
 AX^9+BY^9=1,\qquad X,Y\in k_3^\times.
\]
The first summand has valuation congruent to one modulo nine. It
cannot have negative valuation, because then the second would have
the same valuation, whereas its valuation is zero modulo nine.
It therefore has positive valuation. The second summand is a unit,
so $Y$ is a unit, and reduction forces $Y\equiv1\pmod\alpha$:
the ninth-power map is the identity on $\mathbb F_3^\times$ and
$B\equiv1$. The two displayed depth bounds imply
$v(1-BY^9)=5$, which is incompatible with $1\pmod9$.

The field $k_3$ has no nontrivial cube root of unity, since the
cyclotomic extension has degree two and $[k_3:\mathbb Q_3]=3$.
If $Y^9=B=U^3$, it follows that $Y^3=U$. Reduction gives
$Y\equiv1\pmod\alpha$, contradicting $v(U-1)=2$.
Thus $B$ is not a ninth power. The zero-cusp condition $B=1$
fails for this reason. The one-cusp condition $A=1$ fails by its
valuation, and the infinity-cusp condition $A=B$ fails by the two
valuations modulo nine; the usual minus sign is itself a ninth
power. This treats all three actual cuspidal Kummer images.

For the claimed cubic comparison, set $X=\alpha^2$ and
$t=AX^3$. Then $v(t)=7$. The polynomial $Z^3-(1-t)$ at $Z=1$
has error valuation seven and derivative valuation three, so
$7>2\cdot3$ supplies a root $Y\in k_3^\times$ by Newton's
criterion. Consequently
$t=AX^3$ and $1-t=B(Y/U)^3$; both coordinates are nonzero and
$t\ne1$. This is an actual ordinary point with the prescribed
cubic classes.

The independent complete enumeration in
\texttt{code/check\_a196\_local.py} uses all $729$ elements of
$(\mathbb Z/9)[\alpha]/(\alpha^3-3)=\mathcal O_{k_3}/(\alpha^6)$.
It verifies the identities, the $243$ principal-unit ninth powers,
the absence of integral ordinary solutions and both root exclusions.
The valuation argument above covers nonintegral variables and proves
the claim independently of a finite search heuristic.
\end{proof}
